% ID: SEL-D0037-OB-E01
% ESTADO: validada. S038, cotejada con examen y criterios.
Sean $x,y$ las bolsas A y B. Las disponibilidades dan
$3x+2y\le600$, $x+2y\le400$ y $x,y\ge0$.
Maximizamos los ingresos $I=4x+3y$. Sobre los ejes
se obtienen $(200,0)$ y $(0,200)$. Al restar las ecuaciones
frontera resulta $2x=200$, de donde el cruce es $(100,150)$.
\begin{center}
\setlength{\unitlength}{8mm}
\begin{picture}(8,7)(-1,-1)
\qbezier(0,0)(2.5,0)(5,0)\put(5.2,0){$x$}
\qbezier(0,0)(0,2.5)(0,5)\put(0,5.2){$y$}
\qbezier(0,4)(1,3.5)(2,3)
\qbezier(2,3)(3,1.5)(4,0)
\put(0.1,4.2){$(0,200)$}
\put(2.2,3.15){$(100,150)$}
\put(4,-0.4){\makebox(0,0)[t]{$(200,0)$}}
\end{picture}
\end{center}
En $(0,0)$, $(200,0)$, $(100,150)$ y $(0,200)$, los ingresos
son 0, 800, 850 y 600 euros. La función lineal alcanza su
máximo en el mejor vértice y éste ya tiene coordenadas enteras.
Por tanto se elaboran \textbf{100 bolsas A y 150 bolsas B}
y se ingresan $\boxed{850}$ euros.
